\section{The construction and its analytical content}
\label{sec:overview}

The central object is the even entire function
\begin{equation}
 H_\lambda(z)=\int_0^\infty
 \ee^{\lambda v^2}\Phi(v)\cos(zv)\dd v,
 \qquad
 \Phi(v)=\sum_{m\ge1}
 \bigl(2\pi^2m^4\ee^{9v}-3\pi m^2\ee^{5v}\bigr)
 \ee^{-\pi m^2\ee^{4v}},\quad v\ge0.
 \label{eq:intro-H}
\end{equation}
The source has a smooth even extension to the full line. Its rapid decay
makes every real deformation parameter $\lambda$ admissible. With the
normalization in \eqref{eq:intro-H},
$H_0(z)=\xi((1+\ii z)/2)/8$ and
$\partial_\lambda H_\lambda=-\partial_z^2H_\lambda$
\cite{RodgersTao2020}. We will derive both identities from the source.

The auxiliary construction has three layers. First, the complete cylinder
Hamiltonian supplies a Green weight and a heat semigroup. Second, a
specified scan of the angular damping produces $\Phi$, which is placed
as a spatial amplitude. Third, an explicitly calibrated linear readout
observes the evolved amplitude. Keeping these layers distinct isolates
what the geometry proves from what the preparation and detector supply.

The resulting theorem is simple to state. Write
\begin{equation}
 Q_1(\beta)=\sum_{n\ge1}\ee^{-\beta n^2},\qquad
 q(\tau)=\frac{Q_1(\pi+4\tau)}{Q_1(\pi)}.
 \label{eq:intro-q}
\end{equation}
The normalized angular bank, the prepared field, and the readout are
defined precisely in Sections~\ref{sec:preparation} and~\ref{sec:readout}.

\begin{theorem}[Complete response, overview]
\label{thm:overview}
For the specified source preparation and matched linear readout, the
left-acting heat evolution on the complete cylinder gives
\begin{equation}
 \boxed{\Obs_{\lambda,\tau}(z)
       =q(\tau)\ee^{-\tau z^2}H_\lambda(z).}
 \label{eq:response-overview}
\end{equation}
This holds for every real $\lambda$, every $\tau\ge0$, and every $z\in\C$.
For any fixed finite $M,T,Z,b\ge0$, all fixed finite mixed derivatives
in $\lambda,\tau,z$ are continuous on
\begin{equation}
 |\lambda|\le M,\quad 0\le\tau\le T,\quad
 |\Re z|\le Z,\quad |\Im z|\le b,
 \label{eq:windows}
\end{equation}
and can be passed through the complete source, heat, Fourier, and trace
operations. Proper-time derivatives at zero are right derivatives.
Every bound depends on the specified finite derivative orders and window.
\end{theorem}

We first realize the complete self-adjoint operators, then derive
summable angular tails and integrable radial majorants, and finally
evaluate the response. This order also explains why fixed-frequency
observation, although not bounded on a generic radial $L^2$ space, is
well-defined on the prepared and evolved family.

There are several distinct parameters:
\begin{center}
\begin{tabularx}{\textwidth}{@{}lY@{}}
\toprule
Parameter & Meaning \\
\midrule
$v$ & Dimensionless physical cylinder coordinate. \\
$u$ & External scan coordinate, later imprinted as $u=v$. \\
$\beta$ & Angular damping parameter; the source scan uses $\beta=\pi\ee^{4u}$. \\
$\lambda$ & Newman parameter in the prepared source $\ee^{\lambda v^2}\Phi(v)$. \\
$\tau$ & Nonnegative dimensionless heat or Euclidean proper time. \\
$k$, $z$ & Real radial momentum and its complex continuation. \\
\bottomrule
\end{tabularx}
\end{center}
In particular, increasing $\tau$ is not increasing $\lambda$. Ordinary
unitary laboratory evolution is a third evolution, discussed below.

The parent manuscript provides the common background for the direct and
auxiliary branches \cite{TNDBackground}. The present manuscript develops
the auxiliary branch on its own. Its principal mathematical input is the
complete-spectrum response note \cite{TNDAux}, accepted within its
conditional scope by the branch assessment \cite{TNDAuxReview}, together
with the shared geometry and source analysis
\cite{TNDParent,TNDSharedReview}. The first-jet conditioning program
\cite{TNDAuxConditioning} remains pending; it is separated from the
established results in Section~\ref{sec:outstanding}.
